\section{Diameter of a Binary Tree}
\label{sec:trees-diameter}

\problemstatement{The diameter of a binary tree is the length (number of
edges) of the longest path between any two nodes. The path need not pass
through the root. Given the root, return the diameter.}

\begin{patternbox}
\emph{``Longest path between any two nodes, not necessarily through the
root''} is the classic ``return one thing, track another'' pattern. At
each node the \emph{best path passing through it} is left height $+$ right
height, but what the node must \emph{return} to its parent is only its own
height. Two different quantities: a \textbf{global maximum} for the answer,
a \textbf{return value} for the parent.
\end{patternbox}

\subsection*{Understanding the problem carefully}

Any path has a single highest node -- its ``peak.'' The path through a
given peak descends into the left subtree and the right subtree, so its
length is (height of left subtree) $+$ (height of right subtree) edges.
The overall diameter is the maximum of that quantity over all possible
peaks. Since every node is the peak of exactly one candidate path, we
evaluate this candidate at every node during a height computation.

\begin{ideabox}[Compute Height, Update a Global Diameter Along the Way]
Run the ordinary height recursion. At each node, before returning, update
a global answer with $L+R$ (the through-this-node path length in edges).
Then return $1+\max(L,R)$ to the parent as usual. The node contributes its
own candidate to the global max but hands the parent only what the parent
needs.
\end{ideabox}

\subsection*{Algorithm}

\begin{enumerate}
  \item Keep a global \textit{best} $=0$.
  \item \textsc{Height}(\textit{node}):
  \begin{enumerate}
    \item If null, return $0$.
    \item $L=$\textsc{Height}(left), $R=$\textsc{Height}(right).
    \item \textit{best} $=\max(\textit{best},\ L+R)$.
    \item Return $1+\max(L,R)$.
  \end{enumerate}
  \item Answer is \textit{best}.
\end{enumerate}

\subsection*{Dry run}

\dryrun{Tree: root $1$; left $2$ with children $4,5$; right $3$.}

\begin{itemize}
  \item Height($4$)$=$Height($5$)$=1$. At $2$: $L=1,R=1$, best$=\max(0,2)=2$;
        return $2$.
  \item Height($3$)$=1$. At $1$: $L=2,R=1$, best$=\max(2,3)=3$; return
        $3$.
  \item Diameter $=3$ edges (path $4\!-\!2\!-\!1\!-\!3$).
\end{itemize}

\subsection*{C++ implementation}

\begin{lstlisting}
#include <bits/stdc++.h>
using namespace std;

struct TreeNode {
    int val;
    TreeNode* left;
    TreeNode* right;
    TreeNode(int v) : val(v), left(nullptr), right(nullptr) {}
};

int best;

int height(TreeNode* node) {
    if (node == nullptr) return 0;
    int L = height(node->left);
    int R = height(node->right);
    best = max(best, L + R);      // path through this node, in edges
    return 1 + max(L, R);
}

int diameterOfBinaryTree(TreeNode* root) {
    best = 0;
    height(root);
    return best;
}

int main() {
    TreeNode* root = new TreeNode(1);
    root->left = new TreeNode(2);
    root->right = new TreeNode(3);
    root->left->left = new TreeNode(4);
    root->left->right = new TreeNode(5);

    cout << "Diameter: " << diameterOfBinaryTree(root) << "\n";
    return 0;
}
\end{lstlisting}

\begin{complexitybox}
\textbf{Time Complexity.} $O(n)$ -- one height pass with $O(1)$ extra work
per node. A naive ``compute height twice at every node'' solution would be
$O(n^2)$. \textbf{Space Complexity.} $O(h)$ for the recursion stack.
\end{complexitybox}

\begin{takeawaybox}
``Return one value to the parent, update a separate global for the answer''
is the reusable shape for any tree problem where the optimum can bend
through a node rather than continue straight down. The very next section,
maximum path sum, is the same idea with values instead of heights.
\end{takeawaybox}
